% Part: second-order-logic
% Chapter: sol-and-set-theory
% Section: cardinalities

\documentclass[../../../include/open-logic-section]{subfiles}

\begin{document}

\olfileid{sol}{set}{crd}

\olsection{ગણોનાં કાર્ડિનલ કદ}

\begin{explain}
વ્યાપકક્ષેત્ર સાન્ત કે અનંત, !!{enumerable} કે !!{nonenumerable}
છે એવું વ્યક્ત કરી શકીએ છીએ. તેવી જ રીતે~$\Domain{M}$નો કોઈ ઉપગણ
સાન્ત કે અનંત, !!{enumerable} કે !!{nonenumerable} છે તે ગુણધર્મ
પણ વ્યાખ્યાયિત કરી શકીએ છીએ.
\end{explain}

\begin{prop}
સૂત્ર $\fn{Inf}(X) \ident$
\begin{multline*}
\lexists[u][(\lforall[x][\lforall[y][(\eq[u(x)][u(y)] \lif
      \eq[x][y])]] \land {} \\
\lforall[x][(X(x) \lif X(u(x)))] \land {} \\
\lexists[y][(X(y) \land \lforall[x][(X(x)
      \lif \eq/[y][u(x)])])])]
\end{multline*}
ચલ-નિયુક્તિ~$s$ને અનુલક્ષીને ત્યારે અને માત્ર ત્યારે જ સંતોષાય છે
જ્યારે~$s(X)$ અનંત હોય.
\end{prop}
\sourcecorrection{OLSOL-015}{સ્થિર અંગ્રેજી મૂળના સૂત્રમાં~$u$ને
$X$માં જ મૂલ્યો લેવાની શરત નથી, અને સૂત્રનો બહારનો કૌંસ બંધ થયો
નથી. બે ઘટકોના વ્યાપકક્ષેત્રમાં એક-ઘટક
$X$ લો અને~$u$થી બંને ઘટકો અદલો: મૂળ સૂત્ર સાચું થાય છે, છતાં~$X$
સાન્ત છે. અહીં $X(x) \lif X(u(x))$ ઉમેર્યું છે. ત્યારે~$u$ની
$X$ પરની મર્યાદા એક-એક પરંતુ~$X$ પર વ્યાપ્ત નહીં હોય ત્યારે જ
સૂત્ર સંતોષાય છે; ચૂકી ગયેલો બંધ કૌંસ પણ ઉમેર્યો છે.}

\begin{prop}
સૂત્ર $\fn{Count}(X) \ident $
\begin{multline*}
\lnot \lexists[x][X(x)] \lor {} \\
\lexists[z][\lexists[u][(X(z) \land
    \lforall[x][(X(x) \lif X(u(x)))] \land {} \\ \lforall[Y][((Y(z) \land
      \lforall[x][(Y(x) \lif Y(u(x)))]) \lif X \subseteq Y)])]]
\end{multline*}
ચલ-નિયુક્તિ~$s$ને અનુલક્ષીને ત્યારે અને માત્ર ત્યારે જ સંતોષાય છે
જ્યારે~$s(X)$ !!{enumerable} હોય.
\end{prop}
\sourcecorrection{OLSOL-016}{સ્થિર અંગ્રેજી મૂળ દરેક~$Y$ માટે
$X=Y$ માગે છે; પરંતુ આખું વ્યાપકક્ષેત્ર પોતે~$z$ ધરાવતું અને~$u$
હેઠળ બંધ~$Y$ છે, તેથી મૂળ સૂત્ર કોઈ ઉચિત ઉપગણ~$X$ માટે સાચું થઈ
શકતું નથી. જરૂરી નિષ્કર્ષ $X\subseteq Y$ છે: $z$ની~$u$-કક્ષા
$X$ના બધા ઘટકો સુધી પહોંચે છે. ખાલી ઉપગણ પણ પરિગણનીય ગણાય તે
માટે તેની અલગ વિસંયોજિકા ઉમેરી છે.}

કૅન્ટૉરના પ્રમેયથી આપણે જાણીએ છીએ કે !!{nonenumerable} ગણો છે;
અનંત કદનાં અનંત જુદાં સ્તરો પણ છે. ગણસિદ્ધાંત ગણોનાં કદનું આખું
અંકગણિત વિકસાવે છે અને ગણોને અનંત કાર્ડિનલ સંખ્યાઓ આપે છે. પ્રાકૃતિક
સંખ્યાઓ સાન્ત ગણોનાં કદ માપતી કાર્ડિનલ સંખ્યાઓ તરીકે કામ કરે છે.
!!{denumerable} ગણોનું કદ પહેલું અનંત કાર્ડિનલ કદ છે, જેને~$\aleph_0$
(``આલેફ-નૉટ'' અથવા ``આલેફ-શૂન્ય'') કહે છે. પછીનું અનંત કદ~$\aleph_1$
છે: સૌથી નાનું અપરિગણનીય કાર્ડિનલ કદ. ``$X$નું કદ~$\aleph_0$ છે''
તેને $\fn{Aleph}_0(X) \liff \fn{Inf}(X) \land \fn{Count}(X)$
થી વ્યાખ્યાયિત કરી શકીએ છીએ. $X$નું કદ~$\aleph_1$ ત્યારે અને માત્ર
ત્યારે જ જ્યારે~$X$ પોતે અપરિગણનીય હોય અને તેનો દરેક અપરિગણનીય
ઉપગણ~$X$ સાથે સામ્ય હોય. તેથી આ ગુણધર્મ સૂત્ર
$\fn{Aleph}_1(X) \ident \lnot\fn{Count}(X) \land
\lforall[Y][((Y \subseteq X \land \lnot\fn{Count}(Y))
  \lif \cardeq{X}{Y})]$ દ્વારા વ્યક્ત કરી શકાય છે. $\aleph_2$નું
કદ પણ એ જ રીતે વ્યાખ્યાયિત કરી શકાય છે; વગેરે.
\sourcecorrection{OLSOL-017}{સ્થિર અંગ્રેજી મૂળ કહે છે કે
$\aleph_1$-કદના~$X$ના બધા ઉપગણો સાન્ત અથવા~$\aleph_0$-કદના છે;
પરંતુ~$X$ પોતે પોતાનો ઉપગણ હોવાથી આ અશક્ય છે. અહીં અપરિગણનીય
ઉપગણો~$X$ સાથે સામ્ય હોવાની લઘુત્તમતા-શરત મૂકી છે, અને
\texttt{\textbackslash fn\{Aleph\_1\}}ને અગાઉના
\texttt{\textbackslash fn\{Aleph\}\_0} સાથે સુસંગત
\texttt{\textbackslash fn\{Aleph\}\_1} કર્યું છે.}

એક કદ ખાસ રસનું છે: સાતત્યકનું કાર્ડિનલ કદ. તે~$\Pow{\Nat}$નું
અથવા, સમકક્ષ રીતે, $\Real$નું કદ છે. કોઈ ગણ સાતત્યકના કદનો છે
તે દ્વિતીય-ક્રમ તર્કશાસ્ત્રમાં વ્યક્ત કરી શકાય છે, જોકે તે માટે
થોડું વધુ કામ કરવું પડે છે.

\end{document}
